\chapter{Measures}

\begin{center}\itshape
Every interval has a length. Can every set of real numbers be given a length, so that lengths add up over countably
many disjoint pieces and do not change when a set is shifted? No. Then which sets can, and how many lengths are there?
\end{center}

In 1902 a young French schoolteacher, Henri Lebesgue, defended a thesis with a modest title, \emph{Int\'egrale, longueur,
aire}. In it he asked a question that sounds like a children's riddle: can we give every set of real numbers a length, so that
an interval $[a,b]$ has length $b-a$, shifting a set does not change its length, and the length of a union of countably many
disjoint pieces is the sum of their lengths? He could not answer it. Three years later the Italian Giuseppe Vitali did, and the
answer was no. There are sets that cannot be measured, and the axiom of choice of Chapter~1 is to blame. In 1914 Hausdorff made it
worse, and in 1924 Banach and Tarski made it spectacular: a ball can be cut into finitely many pieces and reassembled into two balls of the
same size. Volume, it turns out, is not defined for every set, and cannot be.

So the question becomes: which sets \emph{can} be measured, and how do we build the measure? The answer, due to Constantin
Carath\'eodory (1914), is a small masterpiece of laziness. Start from something that obviously deserves a length, the boxes. Give
every set the smallest total volume of a countable cover by boxes; this ``outer measure'' is defined for all sets but is not
additive. Then keep only the sets that split every other set additively. They form a $\sigma$-algebra, containing all the sets anyone
will ever meet, and on it the outer measure is a true measure.

Tome~3 built probability on measures and admitted two facts: Lebesgue measure exists, and so do infinite sequences of
independent random variables with prescribed laws. This chapter proves both. The second has a delightful proof: all the coin tosses
of a lifetime, and infinitely many more, are hidden in the binary digits of a single random number. And it ends with Borel's
theorem of 1909, a paradox of a gentler kind: almost every real number is normal, its digits behaving like those of a fair
die, yet nobody has ever proved that a single familiar number, $\sqrt2$, $\pi$ or $e$, is normal.

\section{\texorpdfstring{$\sigma$}{Sigma}-algebras and measures}

\begin{definition}
A \emph{$\sigma$-algebra} on a set $E$ is a family $\mathcal A$ of subsets of $E$ containing $E$ and stable under complements and countable unions (hence under
countable intersections and differences). A \emph{measure} on $(E,\mathcal A)$ is a map $\mu:\mathcal A\to[0,+\infty]$ with $\mu(\varnothing)=0$ that is \emph{countably additive}: $\mu(\bigsqcup_kA_k)=\sum_k\mu(A_k)$
for disjoint $A_k\in\mathcal A$ (the symbol $\bigsqcup$ marks a disjoint union). It is a \emph{probability} if $\mu(E)=1$, \emph{finite} if $\mu(E)<\infty$, and \emph{$\sigma$-finite} if $E$ is a countable
union of sets of finite measure. A set $N$ with $\mu(N)=0$ is \emph{null}, and a property true outside a null set holds \emph{almost everywhere} (a.e.).
\end{definition}

An intersection of $\sigma$-algebras is a $\sigma$-algebra, so every family $\mathcal C$ of subsets generates a smallest $\sigma$-algebra $\sigma(\mathcal C)$ containing it. The
\emph{Borel} $\sigma$-algebra $\mathcal B(X)$ of a topological space is the one generated by the open sets.

\begin{proposition}\label{prop:measureprops}
Let $\mu$ be a measure on $(E,\mathcal A)$ and $A,B,A_k\in\mathcal A$.
\begin{enumerate}[label=(\alph*)]
\item If $A\subseteq B$ then $\mu(A)\le\mu(B)$, and $\mu(B\setminus A)=\mu(B)-\mu(A)$ if $\mu(A)<\infty$.
\item $\mu(\bigcup_kA_k)\le\sum_k\mu(A_k)$ (countable subadditivity).
\item If $A_1\subseteq A_2\subseteq\cdots$, then $\mu(\bigcup_kA_k)=\lim_k\mu(A_k)$ (continuity from below).
\item If $A_1\supseteq A_2\supseteq\cdots$ and $\mu(A_1)<\infty$, then $\mu(\bigcap_kA_k)=\lim_k\mu(A_k)$ (continuity from above).
\item (Borel--Cantelli) If $\sum_k\mu(A_k)<\infty$, then almost every point lies in only finitely many $A_k$.
\end{enumerate}
\end{proposition}
\begin{proof}
(a) $\mu(B)=\mu(A)+\mu(B\setminus A)$. (b) The sets $A'_k=A_k\setminus(A_1\cup\cdots\cup A_{k-1})$ are disjoint, with the same union, and $\mu(A'_k)\le\mu(A_k)$. (c) With $A_0=\varnothing$,
$\bigcup_kA_k=\bigsqcup_k(A_k\setminus A_{k-1})$, so $\mu(\bigcup_kA_k)=\lim_n\sum_{k\le n}\mu(A_k\setminus A_{k-1})=\lim_n\mu(A_n)$. (d) Apply (c) to $A_1\setminus A_k$ and use (a). (e) The set of
points in infinitely many $A_k$ is $\bigcap_m\bigcup_{k\ge m}A_k$, of measure at most $\sum_{k\ge m}\mu(A_k)$ for every $m$, which tends to $0$.
\end{proof}

The Borel sets of $\R^d$ are generated by the \emph{boxes} $]a,b]=\prod_i\,]a_i,b_i]$: a box is the intersection of the open boxes $\prod_i\,]a_i,b_i+\frac1n[$, so it is Borel; and an open
set $U$ is the union of the countably many \emph{dyadic} boxes $\prod_i\,]k_i2^{-n},(k_i+1)2^{-n}]$ contained in it (around $x\in U$, a ball $B(x,r)\subseteq U$ contains the dyadic box of
side $2^{-n}<r/\sqrt d$ that contains $x$).

\section{Uniqueness: Dynkin's theorem}

\begin{definition}
A family $\mathcal P$ of subsets of $E$ is a \emph{$\pi$-system} if it is stable under finite intersections. A family $\mathcal L$ is a \emph{$\lambda$-system} if $E\in\mathcal L$, $B\setminus A\in\mathcal L$ whenever
$A\subseteq B$ are in $\mathcal L$, and $\bigcup_kA_k\in\mathcal L$ for every increasing sequence in $\mathcal L$.
\end{definition}

\begin{theorem}[Dynkin's $\pi$-$\lambda$ theorem]\label{thm:dynkin6}
If a $\lambda$-system $\mathcal L$ contains a $\pi$-system $\mathcal P$, then it contains $\sigma(\mathcal P)$.
\end{theorem}
\begin{proof}
Let $\mathcal D$ be the smallest $\lambda$-system containing $\mathcal P$ (an intersection of $\lambda$-systems is one); $\mathcal D\subseteq\mathcal L$. It suffices to show that $\mathcal D$ is a $\pi$-system:
a $\lambda$-system stable under finite intersections is a $\sigma$-algebra (complements are differences with $E$; finite unions are complements of
intersections of complements; countable unions are increasing unions of finite ones). For $A\in\mathcal D$, the family $\mathcal D_A=\{B:A\cap B\in\mathcal D\}$ is a
$\lambda$-system (since $A\cap(C\setminus B)=(A\cap C)\setminus(A\cap B)$ and intersections with $A$ preserve increasing unions). If $A\in\mathcal P$, then $\mathcal D_A\supseteq\mathcal P$, so $\mathcal D_A\supseteq\mathcal D$:
$A\cap B\in\mathcal D$ for $A\in\mathcal P$, $B\in\mathcal D$. Hence for $B\in\mathcal D$, $\mathcal D_B\supseteq\mathcal P$, so $\mathcal D_B\supseteq\mathcal D$: $\mathcal D$ is a $\pi$-system.
\end{proof}

\begin{theorem}[uniqueness of measures]\label{thm:uniq6}
Let $\mu,\nu$ be measures on $\sigma(\mathcal P)$, where $\mathcal P$ is a $\pi$-system, that agree on $\mathcal P$. If $E=\bigcup_nE_n$ with $E_1\subseteq E_2\subseteq\cdots$ in $\mathcal P$ and $\mu(E_n)<\infty$,
then $\mu=\nu$.
\end{theorem}
\begin{proof}
Fix $n$. The family of $A\in\sigma(\mathcal P)$ with $\mu(A\cap E_n)=\nu(A\cap E_n)$ contains $\mathcal P$ (as $A\cap E_n\in\mathcal P$) and is a $\lambda$-system, by Proposition~\ref{prop:measureprops}(a) and~(c) and the
finiteness of $\mu(E_n)=\nu(E_n)$. By Dynkin's theorem it is all of $\sigma(\mathcal P)$. Let $n\to\infty$ and use continuity from below.
\end{proof}

So two Borel measures on $\R^d$ that agree on boxes, and are finite on bounded boxes, are equal. And a probability on $\R$ is determined
by its distribution function $t\mapsto\mu(]-\infty,t])$, since the half-lines form a $\pi$-system.

\section{Outer measures}

\begin{definition}
An \emph{outer measure} on $E$ is a map $\mu^*:\mathcal P(E)\to[0,+\infty]$ with $\mu^*(\varnothing)=0$, monotone ($A\subseteq B\Rightarrow\mu^*(A)\le\mu^*(B)$) and countably subadditive. A set $A$
is \emph{$\mu^*$-measurable} if it splits every set additively:
\[\mu^*(T)=\mu^*(T\cap A)+\mu^*(T\setminus A)\qquad\text{for every }T\subseteq E.\]
\end{definition}

By subadditivity, only ``$\ge$'' needs checking, and only for $T$ with $\mu^*(T)<\infty$. A set with $\mu^*(N)=0$ is measurable:
$\mu^*(T\cap N)+\mu^*(T\setminus N)\le0+\mu^*(T)$.

\begin{theorem}[Carath\'eodory, 1914]\label{thm:caratheodory}
The $\mu^*$-measurable sets form a $\sigma$-algebra $\mathcal M$, and $\mu^*$ restricted to $\mathcal M$ is a measure. It is \emph{complete}: subsets of null sets are in $\mathcal M$.
\end{theorem}
\begin{proof}
The definition is symmetric in $A$ and $E\setminus A$, so $\mathcal M$ is stable under complements, and $E\in\mathcal M$.

\emph{Finite unions.} Let $A,B\in\mathcal M$ and $T\subseteq E$. Splitting $T$ by $A$, then $T\setminus A$ by $B$:
\[\begin{aligned}\mu^*(T)&=\mu^*(T\cap A)+\mu^*\bigl((T\setminus A)\cap B\bigr)\eqbr+\mu^*\bigl(T\setminus(A\cup B)\bigr)\\&\ge\mu^*\bigl(T\cap(A\cup B)\bigr)+\mu^*\bigl(T\setminus(A\cup B)\bigr),\end{aligned}\]
since $T\cap(A\cup B)=(T\cap A)\cup((T\setminus A)\cap B)$. So $A\cup B\in\mathcal M$, and $\mathcal M$ is an algebra.

\emph{Countable disjoint unions.} Let $A_1,A_2,\ldots\in\mathcal M$ be disjoint, $S_n=A_1\sqcup\cdots\sqcup A_n$ and $S=\bigsqcup_kA_k$. Splitting $T\cap S_n$ by $A_n\in\mathcal M$ gives $\mu^*(T\cap S_n)=\mu^*(T\cap A_n)+\mu^*(T\cap S_{n-1})$,
so by induction $\mu^*(T\cap S_n)=\sum_{k\le n}\mu^*(T\cap A_k)$. As $S_n\in\mathcal M$ and $T\setminus S_n\supseteq T\setminus S$,
\[\begin{aligned}\mu^*(T)&=\mu^*(T\cap S_n)+\mu^*(T\setminus S_n)\eqbr\ge\sum_{k\le n}\mu^*(T\cap A_k)+\mu^*(T\setminus S).\end{aligned}\]
Let $n\to\infty$ and use subadditivity: $\mu^*(T)\ge\sum_k\mu^*(T\cap A_k)+\mu^*(T\setminus S)\ge\mu^*(T\cap S)+\mu^*(T\setminus S)$. So $S\in\mathcal M$. An arbitrary countable union is a disjoint one
of sets $A_k\setminus(A_1\cup\cdots\cup A_{k-1})\in\mathcal M$, so $\mathcal M$ is a $\sigma$-algebra. Taking $T=S$ in the first inequality gives $\mu^*(S)\ge\sum_k\mu^*(A_k)$; the reverse is
subadditivity. So $\mu^*$ is countably additive on $\mathcal M$. Completeness was noted before the theorem.
\end{proof}

\section{Extending a premeasure}

\begin{definition}
A \emph{semiring} on $E$ is a family $\mathcal S$ of subsets containing $\varnothing$, stable under finite intersections, such that $A\setminus B$ is a finite disjoint union of members of
$\mathcal S$ whenever $A,B\in\mathcal S$. A \emph{premeasure} on $\mathcal S$ is a map $\mu_0:\mathcal S\to[0,+\infty]$ with $\mu_0(\varnothing)=0$ that is countably additive on $\mathcal S$: if $S\in\mathcal S$ is the disjoint union
of countably many $S_k\in\mathcal S$, then $\mu_0(S)=\sum_k\mu_0(S_k)$.
\end{definition}

The boxes $]a,b]$ of $\R^d$ (including the empty box) form a semiring: the intersection of two boxes is a box, and if $A=\prod_iA_i$ and $B=\prod_iB_i$ are boxes, each
$A_i$ is the disjoint union of $A_i\cap B_i$ and at most two intervals of the same type, so $A$ is cut into at most $3^d$ disjoint boxes, one of which is $A\cap B$.

\begin{lemma}\label{lem:semiring}
Let $\mathcal S$ be a semiring and $\mu_0:\mathcal S\to[0,\infty]$ finitely additive. (a) For $A,B_1,\dots,B_m\in\mathcal S$, $A\setminus(B_1\cup\cdots\cup B_m)$ is a finite disjoint union of members of $\mathcal S$.
(b) If $S,S_1,\dots,S_m\in\mathcal S$ and $S\subseteq S_1\cup\cdots\cup S_m$, then $\mu_0(S)\le\sum_k\mu_0(S_k)$. If the $S_k$ are disjoint and contained in $S$, then $\sum_k\mu_0(S_k)\le\mu_0(S)$.
\end{lemma}
\begin{proof}
(a) By induction on $m$: if $A\setminus(B_1\cup\cdots\cup B_{m-1})=\bigsqcup_jC_j$, then removing $B_m$ replaces each $C_j$ by $C_j\setminus B_m$, a finite disjoint union of members.
(b) For the second claim, $S\setminus(S_1\cup\cdots\cup S_m)$ is a disjoint union of members by (a), so $S$ is a finite disjoint union of the $S_k$ and these, and additivity
gives the inequality. For the first, $S=\bigsqcup_k\bigl(S\cap S_k\setminus(S_1\cup\cdots\cup S_{k-1})\bigr)$, and by (a) the $k$-th piece is a finite disjoint union of members $R_{kj}\subseteq S_k$. By
additivity and the second claim, $\mu_0(S)=\sum_{k,j}\mu_0(R_{kj})\le\sum_k\mu_0(S_k)$.
\end{proof}

\begin{theorem}[Carath\'eodory's extension theorem]\label{thm:extension}
Let $\mu_0$ be a premeasure on a semiring $\mathcal S$ on $E$. For $A\subseteq E$ let
\[\mu^*(A)=\inf\Bigl\{\sum_k\mu_0(S_k):\ S_k\in\mathcal S,\ A\subseteq\bigcup_kS_k\Bigr\}\]
(with $\inf\varnothing=+\infty$). Then $\mu^*$ is an outer measure, every member of $\mathcal S$ is $\mu^*$-measurable, and $\mu^*=\mu_0$ on $\mathcal S$. Hence $\mu^*$ restricts to a measure on
$\sigma(\mathcal S)$ that extends $\mu_0$; it is the only one if $E$ is a countable increasing union of members of $\mathcal S$ of finite measure.
\end{theorem}
\begin{proof}
\emph{Outer measure.} Monotonicity is clear. For subadditivity, given $A_n$ and $\varepsilon>0$, cover each $A_n$ by members $S_{n,k}$ with $\sum_k\mu_0(S_{n,k})\le\mu^*(A_n)+\varepsilon2^{-n}$; together
they cover $\bigcup_nA_n$, so $\mu^*(\bigcup_nA_n)\le\sum_n\mu^*(A_n)+\varepsilon$.

\emph{$\mu^*=\mu_0$ on $\mathcal S$.} Covering $S$ by itself gives $\mu^*(S)\le\mu_0(S)$. Conversely let $S\subseteq\bigcup_kS_k$. As in Lemma~\ref{lem:semiring}(b), $S$ is the disjoint union of the countably many
members $R_{kj}=$ pieces of $S\cap S_k\setminus(S_1\cup\cdots\cup S_{k-1})$, with $R_{kj}\subseteq S_k$ disjoint for fixed $k$. Countable additivity on $\mathcal S$ and Lemma~\ref{lem:semiring}(b) give
$\mu_0(S)=\sum_k\sum_j\mu_0(R_{kj})\le\sum_k\mu_0(S_k)$. Taking the infimum, $\mu_0(S)\le\mu^*(S)$.

\emph{Members are measurable.} Let $S\in\mathcal S$ and $T$ with $\mu^*(T)<\infty$; cover $T$ by $S_k$ with $\sum_k\mu_0(S_k)\le\mu^*(T)+\varepsilon$. Each $S_k\setminus S$ is a finite disjoint union of members $R_{kj}$,
and $S_k=(S_k\cap S)\sqcup\bigsqcup_jR_{kj}$, so $\mu_0(S_k)=\mu_0(S_k\cap S)+\sum_j\mu_0(R_{kj})$. The $S_k\cap S$ cover $T\cap S$ and the $R_{kj}$ cover $T\setminus S$, so
\[\mu^*(T\cap S)+\mu^*(T\setminus S)\le\sum_k\mu_0(S_k)\le\mu^*(T)+\varepsilon.\]

By Carath\'eodory's theorem, $\mu^*$ is a measure on the $\sigma$-algebra $\mathcal M\supseteq\mathcal S$, hence on $\sigma(\mathcal S)$. Uniqueness is Theorem~\ref{thm:uniq6}, $\mathcal S$ being a $\pi$-system.
\end{proof}

\section{Lebesgue measure}

\begin{figure}[ht]
\centering
\fig{ch06-covering}
\caption{The outer measure of a set (blue) is the infimum of the total volume of countable covers by boxes (orange). Finer covers waste less; for a
Borel set, the infimum is its Lebesgue measure.}
\label{fig:covering}
\end{figure}

\begin{theorem}[Lebesgue measure]\label{thm:lebesgue}
There is a unique measure $\lambda_d$ on $\mathcal B(\R^d)$ with $\lambda_d(\prod_i\,]a_i,b_i])=\prod_i(b_i-a_i)$ for every box. It is translation invariant, $\lambda_d(tA)=t^d\lambda_d(A)$ for $t>0$, and it
is the only translation-invariant measure on $\mathcal B(\R^d)$ that gives the unit cube $]0,1]^d$ measure $1$.
\end{theorem}
\begin{proof}
By Theorem~\ref{thm:extension}, it suffices to show that $\mathrm{vol}(]a,b])=\prod_i(b_i-a_i)$ is a premeasure on the semiring of boxes; uniqueness follows since $\R^d=\bigcup_n\,]-n,n]^d$.

\emph{Finite additivity.} Let a box $Q$ be the disjoint union of boxes $Q_1,\dots,Q_m$. In each coordinate, the endpoints of all these boxes cut the axis into intervals;
their products form a grid of boxes (cells). Each $Q_k$, and $Q$, is the disjoint union of the cells it contains, and its volume is the sum of their volumes:
expanding the product $\prod_i(b_i-a_i)$, each factor being a sum of consecutive cell lengths, gives one term per cell. Each cell of $Q$ lies in exactly one $Q_k$,
so $\mathrm{vol}(Q)=\sum_k\mathrm{vol}(Q_k)$.

\emph{Countable additivity.} Let $Q=]a,b]=\bigsqcup_kQ_k$ with $Q_k=]a^k,b^k]$. For each $N$, Lemma~\ref{lem:semiring}(b) gives $\sum_{k\le N}\mathrm{vol}(Q_k)\le\mathrm{vol}(Q)$, hence $\sum_k\mathrm{vol}(Q_k)\le\mathrm{vol}(Q)$. For the
reverse, let $\varepsilon>0$; we may assume $Q\ne\varnothing$. Since the volume of a box is continuous in its corners, there is $\delta>0$ with $\mathrm{vol}(]a+\delta,b])\ge\mathrm{vol}(Q)-\varepsilon$ (where
$a+\delta$ adds $\delta$ to each coordinate), and for each $k$ there is $\delta_k>0$ with $\mathrm{vol}(]a^k,b^k+\delta_k])\le\mathrm{vol}(Q_k)+\varepsilon2^{-k}$. The closed box $K=[a+\delta,b]$ is compact
(Theorem~\ref{thm:heineborel}) and covered by the open boxes $\prod_i\,]a^k_i,b^k_i+\delta_k[$, so by finitely many, $k\le N$. Then $]a+\delta,b]\subseteq\bigcup_{k\le N}\,]a^k,b^k+\delta_k]$, and by Lemma~\ref{lem:semiring}(b)
\[\mathrm{vol}(Q)-\varepsilon\le\sum_{k\le N}\bigl(\mathrm{vol}(Q_k)+\varepsilon2^{-k}\bigr)\le\sum_k\mathrm{vol}(Q_k)+\varepsilon.\]

\emph{Invariance.} For $c\in\R^d$, $A\mapsto\lambda_d(A+c)$ is a measure on $\mathcal B(\R^d)$ (translation is a homeomorphism, so it maps Borel sets to Borel sets) that agrees with
$\lambda_d$ on boxes, so it equals $\lambda_d$. Likewise $A\mapsto t^{-d}\lambda_d(tA)$.

\emph{Characterisation.} Let $\mu$ be translation invariant with $\mu(]0,1]^d)=1$. The unit cube is the disjoint union of $2^{nd}$ translates of $]0,2^{-n}]^d$, so this cube has
$\mu$-measure $2^{-nd}$. A box with dyadic corners is a finite disjoint union of such cubes, so $\mu$ and $\lambda_d$ agree on it. A general box $]a,b]$ is the increasing union of
boxes $]a^n,b^n]$ with dyadic corners, $a^n$ decreasing to $a$ and $b^n$ increasing to $b$ (strictly above $a$, below $b$), so $\mu$ and $\lambda_d$ agree on it by continuity
from below. By uniqueness, $\mu=\lambda_d$.
\end{proof}

We write $\lambda$ for $\lambda_1$. The sets that are measurable for the outer measure $\lambda_d^*$ of Theorem~\ref{thm:extension} are the \emph{Lebesgue measurable} sets;
they include the Borel sets and all subsets of null sets (Exercise~\ref{ex:completion} describes them exactly). Points, countable sets, lines in $\R^2$
and hyperplanes in $\R^d$ are null (a hyperplane $x_1=c$ is covered by boxes $]c-\varepsilon2^{-k},c]\times\,]-k,k]^{d-1}$ of total volume at most $\varepsilon\sum_k2^{-k}(2k)^{d-1}$, and any hyperplane
is a rotation of this one; Chapter~8 shows that rotations preserve $\lambda_d$). So it does not matter whether boxes are open, closed or half-open.

\begin{proposition}[regularity]\label{prop:regularity}
For every Borel set $A\subseteq\R^d$,
\[\begin{aligned}\lambda_d(A)&=\inf\{\lambda_d(U):U\supseteq A\text{ open}\}\eqbr=\sup\{\lambda_d(K):K\subseteq A\text{ compact}\}.\end{aligned}\]
\end{proposition}
\begin{proof}
\emph{Outer.} We may assume $\lambda_d(A)<\infty$. By the definition of $\lambda_d^*$ there are boxes $]a^k,b^k]$ covering $A$ with $\sum_k\mathrm{vol}\le\lambda_d(A)+\varepsilon$. Enlarge each to an open box
$\prod_i\,]a^k_i,b^k_i+\delta_k[$ of volume at most $\mathrm{vol}(]a^k,b^k])+\varepsilon2^{-k}$. Their union $U$ is open, contains $A$, and $\lambda_d(U)\le\lambda_d(A)+2\varepsilon$.

\emph{Inner.} First let $A$ be bounded, inside a closed box $B$. By the outer part, some open $U\supseteq B\setminus A$ has $\lambda_d(U)\le\lambda_d(B\setminus A)+\varepsilon$. Then $K=B\setminus U$ is closed and
bounded, hence compact, $K\subseteq A$, and $\lambda_d(K)\ge\lambda_d(B)-\lambda_d(U)\ge\lambda_d(A)-\varepsilon$. In general, $A\cap[-n,n]^d$ increases to $A$, and continuity from below finishes.
\end{proof}

So a Borel set is, up to $\varepsilon$, a compact set from inside and an open set from outside; for sets of finite measure this also gives
\emph{Littlewood's first principle}: a Borel set $A$ of finite measure is, up to a set of measure $\varepsilon$, a finite union of boxes (take $U\supseteq A$ open with $\lambda_d(U\setminus A)<\varepsilon$; $U$ is a countable
union of disjoint dyadic boxes (the maximal ones, since two dyadic boxes are nested or disjoint), and by continuity from below finitely many of them have total measure within $\varepsilon$ of $\lambda_d(U)$). Chapter~11 proves regularity for every finite Borel measure on a metric space.

\begin{think}
Can a Borel set $A\subseteq[0,1]$ and its complement both meet every subinterval in positive measure? Regularity says that $A$ looks like a finite union of
intervals, up to $\varepsilon$\dots
\end{think}

Surprisingly, yes (Exercise~\ref{ex:everywhere}). Littlewood's principle approximates in measure, not interval by interval: the error set, of measure
$\varepsilon$, can be spread over every interval.

\section{Lebesgue--Stieltjes measures and laws on \texorpdfstring{$\R$}{R}}

\begin{theorem}\label{thm:stieltjes}
Let $F:\R\to\R$ be non-decreasing and right-continuous. There is a unique Borel measure $\mu_F$ on $\R$ with $\mu_F(]a,b])=F(b)-F(a)$ for all $a<b$. If moreover $F(t)\to0$ as
$t\to-\infty$ and $F(t)\to1$ as $t\to+\infty$, then $\mu_F$ is a probability with distribution function $F$, and every probability on $\mathcal B(\R)$ arises in this way.
\end{theorem}
\begin{proof}
On the semiring of intervals $]a,b]$, $\mu_0(]a,b])=F(b)-F(a)$ is finitely additive, since the increments telescope. For countable additivity, copy the proof of
Theorem~\ref{thm:lebesgue} in dimension one: if $]a,b]=\bigsqcup_k\,]a_k,b_k]$, right-continuity gives $\delta>0$ with $F(a+\delta)\le F(a)+\varepsilon$ and $\delta_k>0$ with $F(b_k+\delta_k)\le F(b_k)+\varepsilon2^{-k}$;
compactness of $[a+\delta,b]$ gives finitely many intervals $]a_k,b_k+\delta_k[$ covering it, and Lemma~\ref{lem:semiring}(b) gives $F(b)-F(a)-\varepsilon\le\sum_k(F(b_k)-F(a_k))+\varepsilon$. Existence and uniqueness
follow from Theorem~\ref{thm:extension} ($\R=\bigcup_n\,]-n,n]$). In the second case, $\mu_F(]-\infty,t])=\lim_n(F(t)-F(-n))=F(t)$ and $\mu_F(\R)=1$, by continuity from below. Conversely the
distribution function $F(t)=\mu(]-\infty,t])$ of a probability is non-decreasing, right-continuous (continuity from above) with the limits $0$ and $1$, and $\mu=\mu_F$ by
uniqueness.
\end{proof}

Lebesgue measure is $\mu_F$ for $F(t)=t$. A jump of $F$ at $a$ is an atom: $\mu_F(\{a\})=F(a)-F(a^-)$ (Exercise~\ref{ex:atoms}). The Cantor staircase of Chapter~2 gives a
continuous probability $\mu_c$ that lives on the null Cantor set; Chapter~10 sorts every law on $\R$ into atoms, a density and such singular parts.

A \emph{random variable} on a probability space $(\Omega,\mathcal A,\mathbb P)$ is a map $X:\Omega\to\R$ with $\{X\in B\}\in\mathcal A$ for every Borel $B$ (a \emph{measurable} function); its \emph{law} is
$B\mapsto\mathbb P(X\in B)$, a probability on $\mathcal B(\R)$. It is enough to check that $\{X\le t\}\in\mathcal A$ for every $t$, because the sets $B$ with $\{X\in B\}\in\mathcal A$ form a $\sigma$-algebra.
For the same reason a sum of random variables is one ($\{X+Y<t\}=\bigcup_{q\in\Q}\{X<q\}\cap\{Y<t-q\}$), and so is the pointwise limit of a sequence ($\{\sup_nX_n\le t\}=\bigcap_n\{X_n\le t\}$, and
$\lim=\inf_m\sup_{n\ge m}$). Chapter~7 develops this.

\begin{proposition}[quantile transform]\label{prop:quantile}
Let $F$ be a distribution function and $G(u)=\inf\{t:F(t)\ge u\}$ for $u\in\,]0,1[$. On $(]0,1[,\mathcal B,\lambda)$, the random variable $G$ has law $\mu_F$.
\end{proposition}
\begin{proof}
$G(u)$ is finite since $F$ has limits $0$ and $1$. By right-continuity the infimum is attained, so $G(u)\le t$ iff $F(t)\ge u$. Hence $\{G\le t\}=\,]0,F(t)]\cap\,]0,1[$, of measure $F(t)$: $G$ is
measurable with distribution function $F$.
\end{proof}

\section{Sets that cannot be measured}

\begin{theorem}[Vitali, 1905]\label{thm:vitali6}
There is a subset of $[0,1]$ that is not Lebesgue measurable. More precisely, no translation-invariant measure on all subsets of $\R$ gives $[0,1]$ a finite
positive measure.
\end{theorem}
\begin{proof}
Call $x,y\in[0,1]$ equivalent if $x-y\in\Q$. By the axiom of choice, let $V\subseteq[0,1]$ contain exactly one point of each class. The translates $V+r$, $r\in\Q\cap[-1,1]$, are
disjoint (if $v+r=v'+r'$, then $v\sim v'$, so $v=v'$ and $r=r'$), contained in $[-1,2]$, and their union contains $[0,1]$ (each $x\in[0,1]$ is equivalent to some $v\in V$, and
$x-v\in\Q\cap[-1,1]$). If $\mu$ were such a measure, defined on $V$, countable additivity and invariance would give $\mu([0,1])\le\sum_r\mu(V)\le\mu([-1,2])\le3\mu([0,1])$: the middle sum of
countably many equal terms is $0$ or $+\infty$, while the ends are finite and positive. Since $\lambda$ on Lebesgue measurable sets is translation invariant (by the
same argument as for Borel sets, translating covers), $V$ is not Lebesgue measurable.
\end{proof}

Choice is essential: Solovay (1970) built a model of set theory without the full axiom of choice (but with countable choices) in which every set of
reals is Lebesgue measurable. In dimension three, Banach and Tarski (1924) cut a ball into finitely many pieces that rigid motions reassemble into two balls of
the same size (five pieces suffice, as Robinson showed in 1947); the pieces cannot be measurable. In dimensions one and two, Banach showed that finitely additive invariant ``measures'' on all sets do
exist, a difference explained by the structure of the groups of motions (amenability, Tome~7, Chapter~8).

There are also Lebesgue measurable sets that are not Borel (Exercise~\ref{ex:nonborel}): the completion adds genuinely new sets, all of them hidden inside null
sets up to a Borel set.

\begin{think}
Is a set of measure zero ``small'' in the sense of Baire? Is a meagre set small in the sense of measure?
\end{think}

Neither: by Theorem~\ref{thm:liouville}, $\R$ is the disjoint union of a null set that is comeagre (the Liouville numbers) and a meagre set of full measure.

\section{Independent sequences from one random number}

How many coins can one toss with a single random number? Write $\omega\in[0,1[$ in binary: if $\omega$ is uniform, its digits are
independent fair coins (Lemma~\ref{lem:digits6}). Now deal the digits like cards to infinitely many players: digits $1,3,5,\ldots$ to the
first, $2,6,10,\ldots$ to the second, and so on. Each player receives infinitely many fair coins, hence a uniform number of her own, and the
players are independent. One random number is infinitely many.

Tome~3 admitted that infinitely many independent random variables with prescribed laws can live on one probability space. Here they are built on
$([0,1[,\mathcal B,\lambda)$ from the binary digits of the point $\omega$. A family of random variables $(X_k)$ is \emph{independent} if for every finite set $J$ of indices and Borel sets $B_k$,
$\mathbb P(X_k\in B_k\text{ for }k\in J)=\prod_{k\in J}\mathbb P(X_k\in B_k)$.

\begin{lemma}\label{lem:digits6}
For $\omega\in[0,1[$ let $D_j(\omega)\in\{0,1\}$ be the $j$-th digit of its binary expansion (the greedy one, Proposition~\ref{prop:cantor}). The $D_j$ are Borel functions, and for any distinct
indices $j_1,\dots,j_N$ and bits $c_1,\dots,c_N$, $\lambda(D_{j_1}=c_1,\dots,D_{j_N}=c_N)=2^{-N}$.
\end{lemma}
\begin{proof}
Let $n\ge\max j_i$. The first $n$ digits of $\omega$ are those of a word $w$ iff $\omega\in[0.w,0.w+2^{-n}[$, an interval of length $2^{-n}$; so each $D_j$ is a step function, Borel. The event in
question is the disjoint union of these intervals over the $2^{n-N}$ words $w$ of length $n$ with the prescribed digits at the positions $j_i$: its measure is $2^{n-N}2^{-n}$.
\end{proof}

\begin{figure}[ht]
\centering
\fig{ch06-digits}
\caption{One uniform number $\omega=0.D_1D_2D_3\ldots$ in base $2$ hides infinitely many. The digits are dealt along the bijection $(k,j)\mapsto2^{k-1}(2j-1)$: the odd positions go to
$U_1$, positions $2,6,10,\ldots$ to $U_2$, positions $4,12,20,\ldots$ to $U_3$, and so on. The rows use disjoint sets of independent digits, so they are independent uniform
numbers.}
\label{fig:digits}
\end{figure}

\begin{theorem}\label{thm:independent}
Let $\mu_1,\mu_2,\ldots$ be probabilities on $\mathcal B(\R)$. On $([0,1[,\mathcal B,\lambda)$ there are independent random variables $X_1,X_2,\ldots$ such that $X_k$ has law $\mu_k$.
\end{theorem}
\begin{proof}
\emph{Splitting.} The map $(k,j)\mapsto2^{k-1}(2j-1)$ is a bijection $\N^*\times\N^*\to\N^*$ (every positive integer is uniquely a power of $2$ times an odd number). Let
\[U_k=\sum_{j\ge1}D_{2^{k-1}(2j-1)}\,2^{-j},\]
a random variable (a pointwise limit of finite sums of Borel functions), with values in $[0,1]$.

\emph{Joint law on dyadic intervals.} Fix $n,m\ge1$ and dyadic intervals $I_k=[i_k2^{-n},(i_k+1)2^{-n}[$, $k\le m$, and let $A_k$ be the event that the first $n$ digits used by $U_k$
spell the binary word of $i_k$. On $A_k$, $U_k$ lies in the closure of $I_k$, and equals its right end only if all later digits of $U_k$ are $1$. Off $A_k$, $U_k$ lies outside $I_k$, unless
it equals the left end $i_k2^{-n}$ through an expansion whose digits are eventually all $1$. So $\{U_k\in I_k\}$ and $A_k$ differ by a subset of the event $Z_k$ that the digits of $U_k$ are
eventually all $1$. For each $J$, the event that digits $D_{2^{k-1}(2j-1)}$ equal $1$ for all $j\ge J$ has measure at most $2^{-M}$ for every $M$ (Lemma~\ref{lem:digits6}, with $M$ of these digits),
so it is null, and $Z_k$, a countable union of such events, is null. Up to a null set, $\{U_k\in I_k\text{ for }k\le m\}$ is $A_1\cap\cdots\cap A_m$, which prescribes $mn$ distinct digits, and Lemma~\ref{lem:digits6} gives
\[\lambda(U_1\in I_1,\dots,U_m\in I_m)=2^{-mn}=\prod_{k\le m}\lambda(I_k).\]
Dyadic intervals of different lengths are finite disjoint unions of dyadic intervals of the smallest length, so the formula holds for all dyadic intervals.

\emph{Extension.} The dyadic intervals of $[0,1[$, with $\varnothing$, form a $\pi$-system generating $\mathcal B([0,1[)$. With $m=1$, the law of $U_k$ agrees with $\lambda$ on it, so $U_k$ is
uniform on $[0,1[$ (Theorem~\ref{thm:uniq6}; note $\lambda(U_k=1)=0$). For independence, fix dyadic $I_2,\dots,I_m$: the Borel $B_1$ with $\lambda(U_1\in B_1,U_2\in I_2,\dots)=\lambda(B_1)\prod_{k\ge2}\lambda(I_k)$
form a $\lambda$-system containing the dyadic intervals, hence all Borel sets (Dynkin). Then fix a Borel $B_1$ and dyadic $I_3,\dots$, and do the same in the second
variable; after $m$ steps, $\lambda(U_k\in B_k,\ k\le m)=\prod_k\lambda(B_k)$ for all Borel $B_k$.

\emph{Laws.} Let $G_k$ be the quantile function of $\mu_k$ (Proposition~\ref{prop:quantile}), and $X_k=G_k(U_k)$ where $U_k\in\,]0,1[$, $X_k=0$ on the null set where $U_k=0$. $G_k$ is non-decreasing, so
$\{G_k\le t\}$ is an interval and $G_k$ is Borel; thus $\{X_k\in B\}=\{U_k\in G_k^{-1}(B)\}$ up to a null set. So $X_k$ has law $\mu_k$, and the $X_k$ inherit independence from the $U_k$.
\end{proof}

This keeps the promise of Tome~3 for sequences of real random variables, (A1) of Tome~3, Chapter~5, and the coins of its Chapter~1: the
digits $D_j$ themselves are independent fair coins. Kolmogorov's extension theorem (Chapter~11) builds, more generally, a sequence with any
consistent family of finite-dimensional joint laws, not only independent ones.

\section{Normal numbers}

Is the digit $7$ as frequent as the digit $3$ in the decimal expansion of $\pi$? Billions of digits have been computed, and the
frequencies look perfectly fair; nobody can prove it, nor even that $7$ appears infinitely often. In 1909 \'Emile Borel turned the
question around: instead of one number, take all of them. Then the answer is yes, for almost every number, in every base at once.

A real number $x\in[0,1[$ is \emph{normal in base $b$} if every block $w$ of $\ell$ digits occurs in its base-$b$ expansion with frequency $b^{-\ell}$: the number of positions $j\le n$
such that digits $j,\dots,j+\ell-1$ spell $w$ is $nb^{-\ell}+o(n)$. It is \emph{absolutely normal} if it is normal in every base $b\ge2$; for $x\in\R$, apply this to its fractional part.

\begin{lemma}\label{lem:fourth}
Let $b\ge2$, $d\in\{0,\dots,b-1\}$, and $N_n(x)$ the number of $j\le n$ with $j$-th base-$b$ digit of $x$ equal to $d$. For $\varepsilon>0$,
\[\lambda\bigl\{x\in[0,1[:\ |N_n(x)-\tfrac nb|\ge n\varepsilon\bigr\}\le\frac3{\varepsilon^4n^2}.\]
\end{lemma}
\begin{proof}
$N_n(x)$ depends only on the first $n$ digits, constant on each of the $b^n$ intervals of length $b^{-n}$ that they define. So the measure in question is $b^{-n}$ times the number
of words $w\in\{0,\dots,b-1\}^n$ with $|\sum_iY_i(w)|\ge n\varepsilon$, where $Y_i(w)=\mathbf 1_{w_i=d}-\frac1b$. This number is at most $\sum_w(\sum_iY_i(w))^4/(n\varepsilon)^4$. Expand the fourth power into terms
$\sum_wY_iY_jY_kY_l$. If one index, say $i$, occurs exactly once, summing over the letter $w_i$ first gives the factor $\sum_{c}(\mathbf 1_{c=d}-\frac1b)=1-1=0$: the term vanishes. The surviving
terms have $i=j=k=l$ ($n$ of them) or two distinct indices each twice ($3n(n-1)$ ordered ones), and each is at most $b^n$ since $|Y_i|\le1$. So the count is at most
$b^n(n+3n(n-1))/(n\varepsilon)^4\le3b^n/(\varepsilon^4n^2)$.
\end{proof}

This is Chebyshev's inequality with a fourth moment, written with words instead of expectations: Tome~3, Chapter~6, did the same for the strong
law of large numbers.

\begin{theorem}[Borel, 1909]\label{thm:borelnormal}
Almost every real number is absolutely normal.
\end{theorem}
\begin{proof}
\emph{Single digits.} With $N_n$ as in Lemma~\ref{lem:fourth} and $\varepsilon=\frac1m$, the bounds $3m^4/n^2$ are summable over $n$, so by Borel--Cantelli, for almost every $x$, $|\frac{N_n(x)}n-\frac1b|<\frac1m$ for
all large $n$. Intersecting over the countably many $m$, $b$ and $d$: for almost every $x$, every digit has frequency $\frac1b$ in every base $b$.

\emph{Blocks.} Fix $b$, $\ell\ge1$, a block $w$ of length $\ell$ and $r\in\{1,\dots,\ell\}$. The positions $j\equiv r\pmod\ell$ cut the expansion of $x$ into consecutive blocks of length $\ell$, which are
the base-$b^\ell$ digits of $y_r=T_r(x)$, the fractional part of $b^{r-1}x$ (shifting a greedy expansion gives the greedy expansion of the shifted number). The map $T_r$ pulls null sets
back to null sets: the preimage of an interval $I\subseteq[0,1[$ is the union of $b^{r-1}$ intervals of length $|I|b^{1-r}$, of total length $|I|$, so a cover of a set $Z$ by intervals of total
length $<\varepsilon$ gives one of $T_r^{-1}(Z)$. By the first step in base $b^\ell$, for almost every $x$, the digit $w$ (read as a number) has frequency $b^{-\ell}$ among the base-$b^\ell$
digits of $y_r$, for every $r$. The occurrences of $w$ at positions $j\le n$ are the sum over $r$ of the occurrences at positions $j\equiv r$; each sum has $\frac n\ell+O(1)$ positions and
$\frac n\ell b^{-\ell}+o(n)$ occurrences, so the total is $nb^{-\ell}+o(n)$. Countably many $b$, $\ell$, $w$, $r$: a countable union of null sets is null.
\end{proof}

\begin{beauty}
Borel's theorem is the strong law of large numbers in disguise, and it was the first: in 1909, before probability had
axioms, Borel used Lebesgue's new measure to say that the digits of a ``random'' number behave like independent dice.
Yet no specific classical constant is known to be normal in any base: not $\sqrt2$, not $\pi$, not $e$. Normal numbers
are everywhere, and we can name almost none of them, apart from constructions such as Champernowne's $0.123456789101112\ldots$ (normal
in base $10$, 1933). And by Proposition~\ref{prop:nonnormal}, in the sense of category, normal numbers are rare.
\end{beauty}

\section{Exercises}

\begin{exercise}[$\star$]\label{ex:atoms}
Let $F$ be non-decreasing and right-continuous. Show that $\mu_F(\{a\})=F(a)-F(a^-)$ and $\mu_F(]a,b[)=F(b^-)-F(a)$. Deduce that $\mu_F$ has no atoms iff $F$ is continuous, and that a
probability on $\R$ has at most countably many atoms.
\end{exercise}

\begin{exercise}[$\star$]\label{ex:fatcantor}
\emph{A fat Cantor set.} From $[0,1]$ remove the open middle interval of length $\frac14$, then from each of the two remaining intervals the open middle interval of
length $\frac1{16}$, and at step $n$ from each of the $2^{n-1}$ remaining intervals the middle interval of length $4^{-n}$. Show that what remains is a closed set with empty
interior and measure $\frac12$. So a nowhere dense set can have positive measure.
\end{exercise}

\begin{exercise}[$\star\star$]\label{ex:cantorsum}
Let $C$ be the Cantor set. Show that $\lambda(C)=0$ and that $C+C=\{x+y:x,y\in C\}$ is the whole interval $[0,2]$. (Use ternary digits, or the compactness of $C\times C$.)
\end{exercise}

\begin{exercise}[$\star\star$]\label{ex:completion}
Show that a set $A\subseteq\R^d$ is Lebesgue measurable iff there are an $F_\sigma$ set $F$ and a $G_\delta$ set $G$ with $F\subseteq A\subseteq G$ and $\lambda_d(G\setminus F)=0$. (So the Lebesgue sets are the
\emph{completion} of the Borel $\sigma$-algebra: Borel sets modified by subsets of null sets.)
\end{exercise}

\begin{exercise}[$\star\star$]\label{ex:steinhaus}
\emph{Steinhaus's theorem.} Let $A\subseteq\R$ be Borel with $\lambda(A)>0$. Show that $A-A=\{x-y:x,y\in A\}$ contains an interval $]-\delta,\delta[$. Deduce that a measurable additive function
$f:\R\to\R$ is linear; compare with Exercise~\ref{ex:hamel}.
\end{exercise}

\begin{exercise}[$\star\star\star$]\label{ex:everywhere}
Construct a Borel set $A\subseteq[0,1]$ such that $0<\lambda(A\cap I)<\lambda(I)$ for every interval $I\subseteq[0,1]$ of positive length. (Put fat Cantor sets inside the gaps of fat Cantor sets, in
every interval with rational ends.)
\end{exercise}

\begin{exercise}[$\star\star\star$]\label{ex:nonborel}
\emph{A Lebesgue measurable set that is not Borel.} Let $c$ be the Cantor staircase (Exercise~\ref{ex:devil}) and $g(x)=x+c(x)$ on $[0,1]$. (a) Show that $g$ is a homeomorphism of $[0,1]$
onto $[0,2]$ and that $\lambda(g(C))=1$. (b) Show that every Lebesgue measurable set $E$ with $\lambda(E)>0$ contains a non-measurable subset. (c) Take a non-measurable
$V\subseteq g(C)$ and show that $g^{-1}(V)$ is Lebesgue measurable but not Borel.
\end{exercise}

\section{Hints}

\begin{hint}{ex:atoms}
$\{a\}=\bigcap_n\,]a-\frac1n,a]$: continuity from above. For the last part, how many atoms can have mass more than $\frac1k$?
\end{hint}
\begin{hint}{ex:fatcantor}
Total length removed: $\sum_n2^{n-1}4^{-n}$. Each remaining interval at step $n$ has length less than $2^{-n}$.
\end{hint}
\begin{hint}{ex:cantorsum}
$\frac x2$ for $x\in C$ has ternary digits $0$ and $1$ only, and every number of $[0,1]$ is a sum of two numbers with ternary digits in $\{0,1\}$ (digit by digit: $0=0+0$, $1=1+0$,
$2=1+1$).
\end{hint}
\begin{hint}{ex:completion}
Use regularity: compact $K_n\subseteq A$ and open $U_n\supseteq A$ with measures within $\frac1n$, first for $A$ bounded. For a Lebesgue set, regularity holds too, with the same proof.
\end{hint}
\begin{hint}{ex:steinhaus}
Take compact $K\subseteq A$ and open $U\supseteq K$ with $\lambda(U)<2\lambda(K)$. If $|t|$ is less than the distance from $K$ to the complement of $U$, then $K+t\subseteq U$, so $K$ and $K+t$ cannot be disjoint.
\end{hint}
\begin{hint}{ex:everywhere}
Enumerate the intervals $I_n$ with rational ends. Build disjoint fat Cantor sets $A_n,B_n$ inside $I_n$, avoiding the (nowhere dense, closed) sets built before, with
$\lambda(A_n)$, $\lambda(B_n)$ small. Let $A=\bigcup_nA_n$.
\end{hint}
\begin{hint}{ex:nonborel}
(a) $g$ is continuous and strictly increasing; on each gap of $C$, $c$ is constant, so $g$ translates the gap. (b) Use Vitali's argument with $E$ in place of $[0,1]$: the sets
$E\cap(V+r)$ cannot all be measurable. (c) A homeomorphism maps Borel sets to Borel sets; $g^{-1}(V)\subseteq C$ is null.
\end{hint}

\section{Solutions}

\begin{solution}{ex:atoms}
\idea Measures of intervals are increments of $F$; atoms are jumps; continuity from above turns ``closed at $a$'' into a left limit.
\writeup $\{a\}=\bigcap_n\,]a-\frac1n,a]$, a decreasing intersection of sets of finite measure, so $\mu_F(\{a\})=\lim_n(F(a)-F(a-\frac1n))=F(a)-F(a^-)$, the left limit existing by monotonicity.
Then $\mu_F(]a,b[)=\mu_F(]a,b])-\mu_F(\{b\})=F(b^-)-F(a)$. So $\mu_F$ has an atom at $a$ iff $F$ jumps at $a$; $F$ is right-continuous, so it is continuous iff it has no jumps. For a
probability, at most $k$ atoms have mass $\ge\frac1k$ (their total is at most $1$), so the atoms form a countable union of finite sets.
\end{solution}

\begin{solution}{ex:fatcantor}
\idea Remove intervals so short that their total length is only $\frac12$, but placed in every remaining interval, so that every interval is eventually cut.
\writeup At step $n$, $2^{n-1}$ intervals of length $4^{-n}$ are removed, a total of $\sum_{n\ge1}2^{n-1}4^{-n}=\frac12\sum_{n\ge1}2^{-n}=\frac12$ (this is possible: by induction each of the $2^{n-1}$ remaining
intervals before step $n$ has length $>4^{-n}$, as checked below). The remaining set $A$ is an intersection of finite unions of closed intervals, so closed, and $\lambda(A)=1-\frac12=\frac12$
by additivity. Let $\ell_n$ be the length of each of the $2^n$ intervals remaining after step $n$. Then $\ell_0=1$ and $\ell_n=\frac12(\ell_{n-1}-4^{-n})<\frac12\ell_{n-1}$, so $\ell_n<2^{-n}$; also
$2^n\ell_n=1-\sum_{k\le n}2^{k-1}4^{-k}>\frac12$, so $\ell_n>2^{-n-1}>4^{-n-1}$. An interval inside $A$ would lie in one remaining interval at each step, so has length $<2^{-n}$ for all $n$: $A$ has empty
interior.
\end{solution}

\begin{solution}{ex:cantorsum}
\idea Halving the Cantor set gives the numbers with ternary digits $0$ and $1$; adding two such numbers digit by digit produces every digit $0$, $1$, $2$ without carries.
\writeup $\lambda(C)\le\lambda(C_n)=(\frac23)^n$ for every $n$ (Proposition~\ref{prop:cantor}), so $\lambda(C)=0$. Let $z\in[0,2]$, so $\frac z2\in[0,1]$ has a ternary expansion $\sum_kt_k3^{-k}$ with $t_k\in\{0,1,2\}$
(greedy, as for binary). Write $t_k=u_k+v_k$ with $u_k,v_k\in\{0,1\}$ ($0=0+0$, $1=1+0$, $2=1+1$). Then $x=\sum_k2u_k3^{-k}$ and $y=\sum_k2v_k3^{-k}$ are in $C$, and $x+y=2\sum_kt_k3^{-k}=z$. Conversely
$C+C\subseteq[0,2]$. So a null set plus itself can be an interval: Steinhaus's theorem (Exercise~\ref{ex:steinhaus}) needs positive measure for $A-A$, and the Cantor set shows that
it is not necessary.
\end{solution}

\begin{solution}{ex:completion}
\idea Regularity sandwiches a set between a compact (or closed) set and an open set with small difference; countably many sandwiches give an $F_\sigma$ and a
$G_\delta$ with null difference.
\writeup The proof of Proposition~\ref{prop:regularity} used only that $A$ is $\lambda_d^*$-measurable (the outer part uses the definition of $\lambda_d^*$, the inner part
complements inside a box), so it applies to Lebesgue sets. If $A$ is Lebesgue measurable and bounded, choose compact $K_n\subseteq A$ and open $U_n\supseteq A$ with
$\lambda_d(U_n\setminus K_n)<\frac1n$; then $F=\bigcup_nK_n$ and $G=\bigcap_nU_n$ work. For unbounded $A$, do this in each $A\cap\,]-m,m]^d$ with error $\frac1n2^{-m}$ and take unions (a countable union of
$F_\sigma$ sets is $F_\sigma$; for $G$, take $U_n=\bigcup_mU_{n,m}$ and $G=\bigcap_nU_n$). Conversely, if $F\subseteq A\subseteq G$ with $\lambda_d(G\setminus F)=0$, then $A\setminus F\subseteq G\setminus F$ is a subset of a null set,
hence $\lambda_d^*$-measurable (completeness, Theorem~\ref{thm:caratheodory}), and $A=F\cup(A\setminus F)$ is measurable.
\end{solution}

\begin{solution}{ex:steinhaus}
\idea A set of positive measure is, by regularity, almost an open set; a small translate of it must overlap itself, since two disjoint copies would not fit in
an open set barely larger.
\writeup By regularity choose compact $K\subseteq A$ with $\lambda(K)>0$, and open $U\supseteq K$ with $\lambda(U)<2\lambda(K)$. The distance $\delta$ from the compact $K$ to the closed set $\R\setminus U$ is positive
(it is attained, Corollary~\ref{cor:extreme}, and $K\cap(\R\setminus U)=\varnothing$). If $|t|<\delta$, then $K+t\subseteq U$. If $K$ and $K+t$ were disjoint, $\lambda(U)\ge\lambda(K)+\lambda(K+t)=2\lambda(K)$ by invariance,
a contradiction. So some $x\in K$ and $y\in K$ have $x=y+t$: $t\in K-K\subseteq A-A$.
For the consequence, let $f$ be additive and Borel. For some $n$, $A=\{|f|\le n\}$ has positive measure (the sets increase to $\R$). By Steinhaus, every $t$ with $|t|<\delta$ is
$x-y$ with $x,y\in A$, so $|f(t)|=|f(x)-f(y)|\le2n$: $f$ is bounded near $0$. Additivity gives $f(qt)=qf(t)$ for rational $q$, so $|f(t)|\le\frac{2n}q$ when $|t|<\frac\delta q$: $f$ is continuous at $0$,
hence everywhere ($f(x+t)-f(x)=f(t)$), and a continuous additive function is $f(x)=f(1)x$ (equal on $\Q$, then by density). The Hamel-basis functions of
Exercise~\ref{ex:hamel} are therefore not measurable.
\end{solution}

\begin{solution}{ex:everywhere}
\idea Every interval with rational ends receives two disjoint fat Cantor sets, one inside $A$ and one outside, each of positive measure. They must avoid what was
built before, which is possible because closed nowhere dense sets leave room, and they must be small, so that the later ones do not fill any interval.
\writeup Let $I_1,I_2,\ldots$ enumerate the open intervals with rational ends in $]0,1[$. A fat Cantor set can be built in any closed interval $J$, scaled (Exercise~\ref{ex:fatcantor}):
it is closed, nowhere dense, and has measure $\frac12|J|$. Build closed nowhere dense sets $A_n,B_n$ by induction. The set $F_{n-1}=\bigcup_{k<n}(A_k\cup B_k)$ is closed and
nowhere dense (a finite union of such sets), so $I_n\setminus F_{n-1}$ is a non-empty open set; choose two disjoint closed intervals $J,J'$ in it, of length at most $2^{-n}$, and let $A_n$, $B_n$ be
fat Cantor sets in $J$ and $J'$. Let $A=\bigcup_nA_n$, a Borel set ($F_\sigma$). Any interval $I$ of positive length contains some $I_n$, so $\lambda(A\cap I)\ge\lambda(A_n)>0$. And $B_n\subseteq I\setminus A$:
$B_n$ is disjoint from $A_k$ for $k\le n$ (by construction) and for $k>n$ (as $A_k$ avoids $F_{k-1}\supseteq B_n$). So $\lambda(I\setminus A)\ge\lambda(B_n)>0$, that is, $\lambda(A\cap I)<\lambda(I)$.
\end{solution}

\begin{solution}{ex:nonborel}
\idea The Cantor staircase squeezes the null Cantor set's complement and stretches $C$ itself to a set of measure $1$. A non-measurable set inside the stretched
copy pulls back to a subset of $C$, measurable because it is null, yet not Borel because homeomorphisms preserve Borel sets.
\writeup (a) $c$ is continuous and non-decreasing with $c(0)=0$, $c(1)=1$, so $g$ is continuous and strictly increasing from $0$ to $2$, hence a continuous bijection $[0,1]\to[0,2]$
(intermediate value theorem), a homeomorphism by Proposition~\ref{prop:compactbasic}(d). On each open gap $]\alpha,\beta[$ of $C$, $c$ is constant, so $g$ maps it onto an interval of the same length. The
gaps have total length $1-\lambda(C)=1$, and their images are disjoint, so $\lambda(g([0,1]\setminus C))=1$ and $\lambda(g(C))=2-1=1$.
(b) Let $V$ be the Vitali set of Theorem~\ref{thm:vitali6}, and $V_r=V+r$ for $r\in\Q$; these are disjoint and cover $\R$ (every $x$ is equivalent to some $v\in V$ after an integer shift; let $V$
contain one point of each class of $\R/\Q$, which is the same as before since every class meets $[0,1]$). If every $E\cap V_r$ were measurable, one of them would have positive
measure, say $\lambda(E\cap V_r)>0$, and then $W=E\cap V_r$, a measurable subset of a translate of $V$, would contain an $F_\sigma$ set $F$ of the same positive measure (Exercise~\ref{ex:completion}), and $W-W\supseteq F-F$ would contain an interval
(Exercise~\ref{ex:steinhaus}). But $W-W\subseteq V-V$ contains no non-zero rational, as two points of $V$ are never equivalent. Contradiction.
(c) By (a) and (b), $g(C)$ contains a non-measurable $V$. The set $g^{-1}(V)\subseteq C$ is a subset of a null set, hence Lebesgue measurable. If it were Borel, then $V=g(g^{-1}(V))$ would
be Borel, since the homeomorphism $g$ maps Borel sets to Borel sets (the sets $B$ with $g(B)$ Borel form a $\sigma$-algebra containing the open sets). But $V$ is not even measurable.
\end{solution}

\section*{Sources and further reading}

\begin{itemize}
\item G.~B.~Folland, \emph{Real Analysis: Modern Techniques and Their Applications}, 2nd ed., Wiley, 1999, Chapter~1. Semirings, premeasures,
Carath\'eodory's theorem and Lebesgue--Stieltjes measures, as here.
\item T.~Tao, \emph{An Introduction to Measure Theory}, American Mathematical Society, 2011, Chapter~1. Lebesgue measure built from elementary
measure and Jordan measure, with many exercises.
\item P.~Billingsley, \emph{Probability and Measure}, 3rd ed., Wiley, 1995, Sections~1--3. Borel's normal number theorem first, then measure theory.
\item H.~Lebesgue, \emph{Int\'egrale, longueur, aire}, Annali di Matematica Pura ed Applicata 7 (1902), 231--359; C.~Carath\'eodory, \emph{Vorlesungen \"uber
reelle Funktionen}, Teubner, 1918.
\item E.~Borel, \emph{Les probabilit\'es d\'enombrables et leurs applications arithm\'etiques}, Rendiconti del Circolo Matematico di Palermo 27 (1909),
247--271. The paper of the papers track: a guide to it comes with this tome.
\item G.~Vitali, \emph{Sul problema della misura dei gruppi di punti di una retta}, Bologna, 1905; R.~M.~Solovay, \emph{A model of set theory in which every
set of reals is Lebesgue measurable}, Annals of Mathematics 92 (1970), 1--56.
\item H.~Steinhaus, \emph{Sur les distances des points dans les ensembles de mesure positive}, Fundamenta Mathematicae 1 (1920), 93--104; D.~G.~Champernowne,
\emph{The construction of decimals normal in the scale of ten}, Journal of the London Mathematical Society 8 (1933), 254--260.
\item \emph{For the stories in the introduction.} Lebesgue's question is in his thesis, cited above, and Vitali's answer in the 1905 pamphlet. F.~Hausdorff, \emph{Grundz\"uge der
Mengenlehre}, Veit, 1914, and S.~Banach and A.~Tarski, ``Sur la d\'ecomposition des ensembles de points en parties respectivement congruentes'', \emph{Fundamenta Math.}
6 (1924), 244--277, for the paradoxes; C.~Carath\'eodory, ``\"Uber das lineare Mass von Punktmengen'', \emph{Nachr.\ Ges.\ Wiss.\ G\"ottingen} (1914), 404--426, is where outer measures come from.
\end{itemize}
